\documentclass[11pt,a4paper]{article}
\usepackage[italian]{babel}
\usepackage[margin=2.5cm]{geometry}
\usepackage{parskip}
\usepackage{amsmath, amssymb, amsthm}
\usepackage{booktabs}
\usepackage{array}
\usepackage{xcolor}
\usepackage{needspace}
\usepackage{enumitem}
\usepackage{listings}
\usepackage{tikz}
\usepackage[most]{tcolorbox}
\usepackage[hidelinks]{hyperref}
\usetikzlibrary{decorations.pathreplacing, shapes.geometric, arrows.meta, positioning}

% Niente vedove/orfane: i paragrafi non lasciano righe isolate a fine/inizio pagina
\widowpenalty=10000
\clubpenalty=10000

% Elenchi più compatti
\setlist{itemsep=2pt, parsep=0pt, topsep=4pt, partopsep=0pt}

% --- Riquadri con bordo nero, senza numero (si spezzano tra due pagine) ---
% Uso: \begin{definizione} ... \end{definizione}
%      \begin{teorema}[Nome] ... \end{teorema}  (il nome è facoltativo)
\tcbset{nota/.style={enhanced jigsaw, breakable, colback=white,
  colframe=black, boxrule=0.5pt, sharp corners}}
\NewTColorBox{definizione}{}{nota,
  before upper={\textbf{Definizione.}\ }}
\NewTColorBox{teorema}{o}{nota,
  before upper={\textbf{Teorema\IfValueT{#1}{ (#1)}.}\ }}
\NewTColorBox{proposizione}{o}{nota,
  before upper={\textbf{Proposizione\IfValueT{#1}{ (#1)}.}\ }}
\NewTColorBox{proprieta}{o}{nota,
  before upper={\textbf{Proprietà\IfValueT{#1}{ (#1)}.}\ }}
\NewTColorBox{assioma}{o}{nota, boxrule=1pt,
  before upper={\textbf{Assioma\IfValueT{#1}{ (#1)}.}\ }}

% --- Ambienti semplici (senza riquadro) ---
% Il titolo sta in cima al blocco, anche se il contenuto inizia con un riquadro o
% con codice; e non resta mai da solo in fondo a una pagina.
% Uso: \begin{esempio}[Titolo facoltativo] ... \end{esempio}
\newcommand{\newrunin}[2]{%
  \NewDocumentEnvironment{#1}{o}{%
    \par\Needspace{4\baselineskip}\addvspace{\medskipamount}\noindent
    \textbf{#2}\IfValueT{##1}{ (##1)}\textbf{.}\ \ignorespaces}%
  {\par\addvspace{\medskipamount}}}
\newrunin{esempio}{Esempio}
\newrunin{esercizio}{Esercizio}
\newrunin{osservazione}{Osservazione}

% V e F nelle tavole di verità
\newcommand{\V}{\textsf{V}}
\newcommand{\F}{\textsf{F}}

% --- Codice C ---
% Uso: \begin{codice} ... \end{codice}      codice C numerato
%      \begin{comando} ... \end{comando}    comandi da terminale
%      \begin{risultato} ... \end{risultato} output di un programma
%      \lstinline|printf("ciao");|           codice dentro al testo
\lstdefinestyle{c}{
  language=C,
  basicstyle=\ttfamily\small,
  keywordstyle=\color{blue}\bfseries,
  commentstyle=\color{gray}\itshape,
  stringstyle=\color{red!70!black},
  numbers=left,
  numberstyle=\tiny\color{gray},
  numbersep=8pt,
  columns=flexible,
  keepspaces=true,
  breaklines=true,
  showstringspaces=false,
  tabsize=4,
  morekeywords={include, define},
  % lettere accentate nei commenti
  literate={à}{{\`a}}1 {è}{{\`e}}1 {é}{{\'e}}1 {ì}{{\`i}}1
           {ò}{{\`o}}1 {ù}{{\`u}}1 {È}{{\`E}}1 {À}{{\`A}}1
}
\lstset{style=c}
\newtcblisting{codice}{listing only, nota, left=6mm,
  listing options={style=c}}
\newtcblisting{comando}{listing only, nota,
  listing options={basicstyle=\ttfamily\small, numbers=none, language={},
    columns=flexible, keepspaces=true, breaklines=true}}
\newtcblisting{risultato}{listing only, enhanced jigsaw, breakable,
  colback=black!5, colframe=black, boxrule=0.5pt, sharp corners,
  listing options={basicstyle=\ttfamily\small, numbers=none, language={},
    columns=flexible, keepspaces=true, breaklines=true}}

% --- Diagrammi di flusso (simboli standard) ---
\tikzset{
  terminale/.style = {ellipse, draw, minimum width=2.5cm, minimum height=0.9cm},
  io/.style        = {trapezium, trapezium left angle=70, trapezium right angle=110,
                      draw, minimum width=2.5cm, minimum height=0.9cm},
  processo/.style  = {rectangle, draw, minimum width=2.5cm, minimum height=0.9cm},
  decisione/.style = {diamond, draw, aspect=2, minimum width=2.5cm},
  freccia/.style   = {thick, -{Stealth}}
}

\title{Esempi di funzioni e curve parametriche}
\author{Marco Cerbella}
\date{Lezione 3 -- 1 ottobre 2026}

\begin{document}
\maketitle

\section{Alcuni esempi di funzioni}

\begin{esempio}[$f_1(x) = x^2$]
$f_1 : \mathbb{R} \to \mathbb{R}$, $f_1(x) = x^2$. L'immagine è $f_1(\mathbb{R}) = [0, +\infty)$: la funzione è limitata inferiormente ma non superiormente. Il grafico
\[
\mathrm{Gr}(f_1) = \{(x, x^2) \mid x \in \mathbb{R}\}
\]
è una parabola con vertice nell'origine e asse coincidente con l'asse $y$.
\end{esempio}

\begin{esempio}[$f_2(x) = 1/x$]
$f_2 : \mathbb{R} \setminus \{0\} \to \mathbb{R}$, $f_2(x) = \dfrac1x$. Il grafico
\[
\mathrm{Gr}(f_2) = \left\{\left(x, \tfrac1x\right) \mid x \in \mathbb{R},\ x \neq 0\right\}
\]
è un'iperbole equilatera.
\end{esempio}

\begin{center}
\begin{tikzpicture}[scale=0.9]
  % parabola
  \draw[->] (-2,0) -- (2,0) node[right] {$x$};
  \draw[->] (0,-0.5) -- (0,3.2) node[above] {$y$};
  \draw[thick, domain=-1.7:1.7, smooth, variable=\x] plot ({\x},{\x*\x});
  \node at (0,-1) {$f_1(x) = x^2$};
  % iperbole
  \begin{scope}[xshift=6cm]
    \draw[->] (-2.5,0) -- (2.5,0) node[right] {$x$};
    \draw[->] (0,-2.5) -- (0,2.5) node[above] {$y$};
    \draw[thick, domain=0.4:2.4, smooth, variable=\x] plot ({\x},{1/\x});
    \draw[thick, domain=-2.4:-0.4, smooth, variable=\x] plot ({\x},{1/\x});
    \node at (0,-3) {$f_2(x) = 1/x$};
  \end{scope}
\end{tikzpicture}
\end{center}

\section{Funzioni definite a tratti}

\begin{esempio}[Valore assoluto]
\[
f_3(x) = |x| = \begin{cases} x & \text{se } x \geq 0, \\ -x & \text{se } x < 0. \end{cases}
\]
Per $x \geq 0$ il grafico coincide con quello di $g(x) = x$ (bisettrice del I e III quadrante); per $x < 0$ con quello di $h(x) = -x$ (bisettrice del II e IV quadrante). Il grafico di $|x|$ è quindi formato dalle due semirette uscenti dall'origine nel I e II quadrante.
\end{esempio}

\begin{esempio}
$f_4 : [0, 3] \to \mathbb{R}$,
\[
f_4(x) = \begin{cases}
3x & \text{se } 0 \leq x \leq 1, \\
4 - x & \text{se } 1 < x \leq 2, \\
x - 1 & \text{se } 2 < x \leq 3.
\end{cases}
\]
Il grafico è formato da tre segmenti: da $(0, 0)$ a $(1, 3)$; da $(1, 3)$ (escluso) a $(2, 2)$; da $(2, 1)$ (escluso) a $(3, 2)$.
\end{esempio}

\begin{center}
\begin{tikzpicture}[scale=1]
  \draw[->] (-2,0) -- (2,0) node[right] {$x$};
  \draw[->] (0,-0.3) -- (0,2.2) node[above] {$y$};
  \draw[thick] (-1.8,1.8) -- (0,0) -- (1.8,1.8);
  \node at (0,-0.9) {$|x|$};
  \begin{scope}[xshift=6cm]
    \draw[->] (-0.5,0) -- (3.7,0) node[right] {$x$};
    \draw[->] (0,-0.3) -- (0,3.5) node[above] {$y$};
    \foreach \x in {1,2,3} \draw (\x,0.05) -- (\x,-0.05) node[below] {$\x$};
    \draw[thick] (0,0) -- (1,3);
    \draw[thick] (1,3) -- (2,2);
    \draw[thick] (2,1) -- (3,2);
    \fill (0,0) circle (1.5pt); \fill (1,3) circle (1.5pt); \fill (2,2) circle (1.5pt); \fill (3,2) circle (1.5pt);
    \draw[fill=white] (2,1) circle (1.5pt);
    \node at (1.5,-0.9) {$f_4(x)$};
  \end{scope}
\end{tikzpicture}
\end{center}

\section{Funzioni a valori vettoriali}

\begin{definizione}
Si dice \emph{funzione a valori vettoriali} una applicazione che ha come codominio un sottoinsieme di $\mathbb{R}^2$ oppure $\mathbb{R}^3$, più in generale di $\mathbb{R}^n$. Un esempio sono le curve parametriche:
\[
f : D \subseteq \mathbb{R} \to \mathbb{R}^2, \qquad t \mapsto f(t).
\]
\end{definizione}

\section{Curve parametriche}

\begin{definizione}
Una \emph{curva parametrica piana} è una funzione $c : [a, b] \to \mathbb{R}^2$,
\[
c(t) = \big(x(t), y(t)\big);
\]
una \emph{curva parametrica nello spazio} è una funzione $c : [a, b] \to \mathbb{R}^3$.
\end{definizione}

\begin{definizione}
Il \emph{supporto} (o \emph{sostegno}) di una curva parametrica è l'insieme di tutti i punti $P = c(t)$ al variare del parametro $t \in [a, b]$.
\end{definizione}

\begin{esempio}[La retta in forma parametrica]
Data una retta $r$ nel piano si scelgono due punti $P_1$ e $P_2$ su di essa. Si individua così un vettore $\vec v = \overrightarrow{P_1P_2}$, detto \emph{vettore direttore} di $r$. Una condizione necessaria e sufficiente affinché un generico punto $P = (x, y)$ appartenga a $r$ è che il vettore $\overrightarrow{P_1P}$ sia parallelo a $\vec v$. Quindi, posto $P_1 = (a, b)$ e $\vec v = (v_1, v_2)$, si ottengono le \emph{equazioni parametriche della retta}:
\[
\begin{cases} x(t) = a + t v_1 \\ y(t) = b + t v_2 \end{cases} \quad t \in \mathbb{R}.
\]
\end{esempio}

\end{document}
